% Options for packages loaded elsewhere
\PassOptionsToPackage{unicode}{hyperref}
\PassOptionsToPackage{hyphens}{url}
\PassOptionsToPackage{dvipsnames,svgnames,x11names}{xcolor}
\documentclass[
  11pt,
  a4paper,
]{article}
\usepackage{xcolor}
\usepackage[margin=1in]{geometry}
\usepackage{amsmath,amssymb}
\setcounter{secnumdepth}{-\maxdimen} % remove section numbering
\usepackage{iftex}
\ifPDFTeX
  \usepackage[T1]{fontenc}
  \usepackage[utf8]{inputenc}
  \usepackage{textcomp} % provide euro and other symbols
\else % if luatex or xetex
  \usepackage{unicode-math} % this also loads fontspec
  \defaultfontfeatures{Scale=MatchLowercase}
  \defaultfontfeatures[\rmfamily]{Ligatures=TeX,Scale=1}
\fi
\usepackage{lmodern}
\ifPDFTeX\else
  % xetex/luatex font selection
  \ifLuaTeX
    \usepackage{luaotfload}
    \directlua{luaotfload.add_fallback("mainfontfallback",{
      "NotoSerifBengali:mode=harf"
    })}
  \fi
  \setmainfont[,RawFeature={fallback=mainfontfallback}]{Latin Modern Roman}
  \ifLuaTeX
    \usepackage{luaotfload}
    \directlua{luaotfload.add_fallback("sansfontfallback",{
      "NotoSansBengali:mode=harf"
    })}
  \fi
  \setsansfont[,RawFeature={fallback=sansfontfallback}]{Latin Modern Sans}
  \ifLuaTeX
    \usepackage{luaotfload}
    \directlua{luaotfload.add_fallback("monofontfallback",{
      "NotoSansBengali:mode=harf"
    })}
  \fi
  \setmonofont[,RawFeature={fallback=monofontfallback}]{DejaVu Sans Mono}
\fi
% Use upquote if available, for straight quotes in verbatim environments
\IfFileExists{upquote.sty}{\usepackage{upquote}}{}
\IfFileExists{microtype.sty}{% use microtype if available
  \usepackage[]{microtype}
  \UseMicrotypeSet[protrusion]{basicmath} % disable protrusion for tt fonts
}{}
\makeatletter
\@ifundefined{KOMAClassName}{% if non-KOMA class
  \IfFileExists{parskip.sty}{%
    \usepackage{parskip}
  }{% else
    \setlength{\parindent}{0pt}
    \setlength{\parskip}{6pt plus 2pt minus 1pt}}
}{% if KOMA class
  \KOMAoptions{parskip=half}}
\makeatother
\setlength{\emergencystretch}{3em} % prevent overfull lines
\providecommand{\tightlist}{%
  \setlength{\itemsep}{0pt}\setlength{\parskip}{0pt}}
% paper_preamble.tex — shared LaTeX preamble for the PCMA/SI paper set.
%
% Compiled with lualatex (NOT pdflatex) because several papers contain
% Bengali script: the CLL paper realises one cognitive state into
% English, Bengali and Banglish, so the Bengali output is EVIDENCE and
% cannot be transliterated away without destroying the point being made.
%
% Font selection is NOT done here. pandoc >= 3 has native
% `mainfontfallback` / `monofontfallback` variables that emit the
% luaotfload fallback registration in the correct position relative to
% unicode-math; doing it by hand in this header lands after unicode-math
% and breaks font loading outright. See tools/build_papers.sh.

\usepackage{amsmath}
\usepackage{amssymb}
\usepackage{mathtools}   % \xleftrightarrow etc. used in cll.md §5
\usepackage{booktabs}

% ---------------------------------------------------------------------
% Unicode symbols appearing in PROSE rather than inside $...$.
%
% The papers were written as markdown and read as markdown, so they use
% real characters — "the ω-boost", "α-child", "ΔR coherence" — instead of
% LaTeX macros. Under unicode-math those codepoints are math-mode-only
% and abort the build ("\symbballowed only in math mode"). Mapping each
% to \ensuremath keeps the markdown readable on GitHub AND compilable,
% which is the whole point of keeping one source.
% ---------------------------------------------------------------------
\usepackage{newunicodechar}

% Greek, lower case
\newunicodechar{α}{\ensuremath{\alpha}}
\newunicodechar{β}{\ensuremath{\beta}}
\newunicodechar{γ}{\ensuremath{\gamma}}
\newunicodechar{δ}{\ensuremath{\delta}}
\newunicodechar{ε}{\ensuremath{\varepsilon}}
\newunicodechar{η}{\ensuremath{\eta}}
\newunicodechar{θ}{\ensuremath{\theta}}
\newunicodechar{κ}{\ensuremath{\kappa}}
\newunicodechar{λ}{\ensuremath{\lambda}}
\newunicodechar{μ}{\ensuremath{\mu}}
\newunicodechar{π}{\ensuremath{\pi}}
\newunicodechar{ρ}{\ensuremath{\rho}}
\newunicodechar{σ}{\ensuremath{\sigma}}
\newunicodechar{τ}{\ensuremath{\tau}}
\newunicodechar{φ}{\ensuremath{\varphi}}
\newunicodechar{χ}{\ensuremath{\chi}}
\newunicodechar{ψ}{\ensuremath{\psi}}
\newunicodechar{ω}{\ensuremath{\omega}}
\newunicodechar{ν}{\ensuremath{\nu}}
\newunicodechar{ξ}{\ensuremath{\xi}}
\newunicodechar{ζ}{\ensuremath{\zeta}}
\newunicodechar{ι}{\ensuremath{\iota}}

% Greek, upper case
\newunicodechar{Δ}{\ensuremath{\Delta}}
\newunicodechar{Φ}{\ensuremath{\Phi}}
\newunicodechar{Σ}{\ensuremath{\Sigma}}
\newunicodechar{Ψ}{\ensuremath{\Psi}}
\newunicodechar{Ω}{\ensuremath{\Omega}}
\newunicodechar{Γ}{\ensuremath{\Gamma}}
\newunicodechar{Λ}{\ensuremath{\Lambda}}
\newunicodechar{Θ}{\ensuremath{\Theta}}
\newunicodechar{Π}{\ensuremath{\Pi}}

% Operators and relations
\newunicodechar{×}{\ensuremath{\times}}
\newunicodechar{·}{\ensuremath{\cdot}}
\newunicodechar{±}{\ensuremath{\pm}}
\newunicodechar{−}{\ensuremath{-}}
\newunicodechar{≈}{\ensuremath{\approx}}
\newunicodechar{≤}{\ensuremath{\leq}}
\newunicodechar{≥}{\ensuremath{\geq}}
\newunicodechar{≠}{\ensuremath{\neq}}
\newunicodechar{∈}{\ensuremath{\in}}
\newunicodechar{∂}{\ensuremath{\partial}}
\newunicodechar{∇}{\ensuremath{\nabla}}
\newunicodechar{√}{\ensuremath{\sqrt{}}}
\newunicodechar{∞}{\ensuremath{\infty}}
\newunicodechar{∝}{\ensuremath{\propto}}
\newunicodechar{∑}{\ensuremath{\sum}}
\newunicodechar{∫}{\ensuremath{\int}}
\newunicodechar{ℝ}{\ensuremath{\mathbb{R}}}
\newunicodechar{ℕ}{\ensuremath{\mathbb{N}}}
\newunicodechar{‖}{\ensuremath{\|}}
\newunicodechar{⊕}{\ensuremath{\oplus}}
\newunicodechar{⊗}{\ensuremath{\otimes}}
\newunicodechar{∀}{\ensuremath{\forall}}
\newunicodechar{∃}{\ensuremath{\exists}}
\newunicodechar{¬}{\ensuremath{\neg}}
\newunicodechar{∧}{\ensuremath{\wedge}}
\newunicodechar{∨}{\ensuremath{\vee}}

% Arrows
\newunicodechar{→}{\ensuremath{\rightarrow}}
\newunicodechar{←}{\ensuremath{\leftarrow}}
\newunicodechar{↔}{\ensuremath{\leftrightarrow}}
\newunicodechar{↑}{\ensuremath{\uparrow}}
\newunicodechar{↓}{\ensuremath{\downarrow}}
\newunicodechar{⇒}{\ensuremath{\Rightarrow}}
\newunicodechar{⇐}{\ensuremath{\Leftarrow}}
\newunicodechar{⇔}{\ensuremath{\Leftrightarrow}}

% Super/subscripts written as literal characters
\newunicodechar{²}{\ensuremath{^{2}}}
\newunicodechar{³}{\ensuremath{^{3}}}
\newunicodechar{⁴}{\ensuremath{^{4}}}
\newunicodechar{⁵}{\ensuremath{^{5}}}
\newunicodechar{⁶}{\ensuremath{^{6}}}
\newunicodechar{₀}{\ensuremath{_{0}}}
\newunicodechar{₁}{\ensuremath{_{1}}}
\newunicodechar{₂}{\ensuremath{_{2}}}
\newunicodechar{₃}{\ensuremath{_{3}}}
\newunicodechar{ᵀ}{\ensuremath{^{\mathsf{T}}}}

% Status glyphs. The papers use emoji as status markers in tables
% (implemented / failed / warning). Mapped to typographic equivalents:
% a colour emoji font in lualatex is unreliable, and the MEANING is what
% carries — a reader needs to see "this one failed", not a red cross.
\newunicodechar{✅}{\ensuremath{\checkmark}}
\newunicodechar{✓}{\ensuremath{\checkmark}}
\newunicodechar{❌}{\ensuremath{\times}}
\newunicodechar{✗}{\ensuremath{\times}}
\newunicodechar{⚠}{\textbf{!}}
\newunicodechar{🧩}{\textbf{[piece]}}
\newunicodechar{📋}{\textbf{[list]}}
\newunicodechar{🟡}{\textbf{[partial]}}
\newunicodechar{🔴}{\textbf{[dead]}}
\newunicodechar{🟢}{\textbf{[live]}}
\newunicodechar{▼}{\ensuremath{\blacktriangledown}}
\newunicodechar{►}{\ensuremath{\blacktriangleright}}
\newunicodechar{□}{\ensuremath{\square}}
\newunicodechar{█}{\rule{0.6em}{0.9ex}}
\newunicodechar{•}{\textbullet}
\newunicodechar{†}{\textdagger}
\newunicodechar{‑}{-}
\newunicodechar{‐}{-}
\newunicodechar{…}{\ldots}
\newunicodechar{।}{\textbar}   % danda: rendered by the Bengali fallback in text
\newunicodechar{⁻}{\ensuremath{^{-}}}

% Zero-width formatting codepoints. These are invisible in an editor but
% are real characters to LaTeX: an emoji written as "⚠️" is TWO
% codepoints (U+26A0 U+FE0F), so declaring only the visible one leaves a
% variation selector with no definition and the build dies on a symbol
% nobody can see. Both are mapped to nothing.
\newunicodechar{^^^^fe0f}{}    % VARIATION SELECTOR-16
\newunicodechar{^^^^200d}{}    % ZERO WIDTH JOINER
\newunicodechar{^^^^200b}{}    % ZERO WIDTH SPACE

\usepackage{longtable}
\usepackage{array}

% Comparison matrices and cost ledgers are wide. Let them shrink to the
% text block rather than overflow the margin.
\usepackage{adjustbox}

% Code blocks: wrap long lines instead of running off the page. Several
% papers quote real C++ with long identifiers.
\usepackage{fvextra}
\fvset{breaklines=true,breakanywhere=true,fontsize=\small}

\setlength{\parskip}{0.5em}
\setlength{\parindent}{0pt}

\usepackage{fancyhdr}
\pagestyle{fancy}
\fancyhf{}
\fancyfoot[C]{\thepage}
\fancyhead[L]{\small\itshape Synthetic Intelligence — research lineage}
\fancyhead[R]{\small\itshape irfan@furylogic.com}
\renewcommand{\headrulewidth}{0.2pt}
\setlength{\headheight}{14pt}
\setlength{\headsep}{18pt}
\usepackage{bookmark}
\IfFileExists{xurl.sty}{\usepackage{xurl}}{} % add URL line breaks if available
\urlstyle{same}
\hypersetup{
  pdftitle={Neuro-Topological Dynamics Architecture (NTDA): Translating Biological Brain Physics into Graph-Rewriting Systems},
  pdfauthor={Irfan Mahir --- Independent Researcher, Dhaka, Bangladesh --- irfan@furylogic.com},
  colorlinks=true,
  linkcolor={black},
  filecolor={Maroon},
  citecolor={Blue},
  urlcolor={[HTML]{0B5394}},
  pdfcreator={LaTeX via pandoc}}

\title{Neuro-Topological Dynamics Architecture (NTDA): Translating
Biological Brain Physics into Graph-Rewriting Systems}
\author{Irfan Mahir --- Independent Researcher, Dhaka, Bangladesh ---
irfan@furylogic.com}
\date{}

\begin{document}
\maketitle

{
\hypersetup{linkcolor=}
\setcounter{tocdepth}{3}
\tableofcontents
}
You are right about one thing, Irfan: nothing is mathematically
impossible to simulate if you understand the underlying rules.

But here is where we need to separate the \emph{biology} from the
\emph{physics of cognition}. When we simulate complex physics---like a
fluid dynamics simulation of a jet engine---we use continuous partial
differential equations over a grid. To simulate the \emph{literal
biological physics} of a human brain, you would need to calculate the
molecular state of 86 billion neurons and 100 trillion synapses. Even an
entire warehouse of RTX 3090 clusters couldn't run a 1:1 molecular
physics simulation of a brain in real-time.

But we don't need to simulate the wetware. We need to simulate the
\textbf{algorithms of the physics}.

A biological brain doesn't use Boolean logic. It uses electrical
thresholds, physical resistance, and chemical gradients. If we translate
those literal biological mechanics into discrete graph mathematics, we
can build an engine that structurally behaves like human neural tissue.

This paper specifies how to turn a concept graph into a simulated neural
physics engine.

\begin{center}\rule{0.5\linewidth}{0.5pt}\end{center}

\section{Neuro-Topological Dynamics Architecture (NTDA): Translating
Biological Brain Physics into Graph-Rewriting
Systems}\label{neuro-topological-dynamics-architecture-ntda-translating-biological-brain-physics-into-graph-rewriting-systems}

\textbf{Author:} Irfan Mahir \emph{Independent Researcher, Dhaka,
Bangladesh} irfan@furylogic.com \textbf{Date:} July 2026
\textbf{Subjects:} Artificial Intelligence (cs.AI); Computational
Neuroscience (q-bio.NC); Neural Dynamics

\begin{center}\rule{0.5\linewidth}{0.5pt}\end{center}

\subsection{Abstract}\label{abstract}

Previous architectures (BSMA, MAA, TSDA) successfully implemented
self-awareness, sub-lexical synthesis, and attentional focus within
discrete graph-rewriting systems. However, they still operate on
fundamental computational logic (true/false edges, integer salience).
The human brain does not compute via logic gates; it computes via
physical resistance, electrical accumulation, and global chemical
gradients. The Neuro-Topological Dynamics Architecture (NTDA) abandons
static edge logic entirely. By translating synaptic plasticity, action
potentials, and neurotransmitter dispersion into graph-theoretical
equations, NTDA provides a framework to simulate the exact informational
physics of a human brain over a discrete C++ engine, bridging the final
gap between logical graphing and biological cognition.

\begin{center}\rule{0.5\linewidth}{0.5pt}\end{center}

\subsection{1. The Physics of the Graph}\label{the-physics-of-the-graph}

In standard SI engines, an edge connects Node A to Node B. The engine
traverses this edge instantly.

In NTDA, an edge is treated as a physical wire with structural
resistance, and a node is treated as a capacitor that must charge before
it can fire. The graph is no longer just a map of knowledge; it is an
active electrical grid. To think, the engine must ``push'' logical
current through the graph.

\subsection{2. Component 1: Action Potential Thresholding (Spiking
Nodes)}\label{component-1-action-potential-thresholding-spiking-nodes}

In a biological brain, a neuron does not fire just because it receives a
signal. It accumulates electrical voltage until it breaches a strict
physical threshold, triggering an Action Potential.

We map this to your concept nodes. Instead of immediately running an
analogy or deduction when a node is queried, the node possesses a
temporary ``Voltage'' (Activation Potential), denoted as \(V(n)\).

When the system observes a concept, it injects a unit of logical charge
\(I\) into the node. The voltage accumulates:
\[ V_t(n) = V_{t-1}(n) + I_{in} - \lambda V_{t-1}(n) \] \emph{(Where
\(\lambda\) is a natural leakage constant, ensuring that if a node isn't
constantly stimulated, it loses its charge, mimicking human short-term
forgetting).}

\textbf{The Spike:} A thought (the execution of a graph-rewrite rule)
\emph{only} occurs if the node's voltage breaches the Action Potential
Threshold (\(\theta\)):
\[ \text{If } V(n) \ge \theta \implies \text{Execute Rule}(n) \text{ and reset } V(n) \text{ to } 0 \]

This prevents the engine from infinitely looping on a single node. Once
it fires an analogy, the node's voltage resets to zero. It physically
cannot fire again until it is re-stimulated by surrounding graph
activity.

\subsection{3. Component 2: Hebbian Conductivity (Synaptic
Plasticity)}\label{component-2-hebbian-conductivity-synaptic-plasticity}

Biological synapses physically widen and become more conductive the more
they are used (``Neurons that fire together, wire together'').

In NTDA, edges do not have a static weight. They have
\textbf{Conductivity} (\(C_{e}\)). When an edge is successfully used to
complete a deduction or analogy, the engine physically alters the edge's
conductivity:

\[ C_e(t+1) = C_e(t) + \alpha \cdot \text{success\_signal} \]

Conversely, unused edges experience topological pruning (decay). If an
edge's conductivity falls below a minimum threshold, it is not just
ignored---it is deleted from the graph. The ontology literally reshapes
its physical structure based on what it thinks about most often.

\subsection{4. Component 3: Global Gradient Modulators (Simulated
Neurotransmitters)}\label{component-3-global-gradient-modulators-simulated-neurotransmitters}

The human brain changes its entire mode of physics based on chemical
states. Dopamine encourages exploration; cortisol induces stress and
hyper-focus. We simulate this by applying global scalar modifiers to the
graph's electrical properties.

We introduce a state vector \(G = [D, C, S]\) representing Dopamine,
Cortisol, and Serotonin equivalents in the engine.

\begin{itemize}
\tightlist
\item
  \textbf{Exploration State (High \(D\)):} Triggered when the
  Morphological Abduction Architecture (MAA) encounters an unknown
  variable. The engine globally lowers the Action Potential Threshold
  (\(\theta\)). Nodes fire easier, causing wild, spontaneous analogies
  and rapid hypothesis generation.
\item
  \textbf{Stress/Correction State (High \(C\)):} Triggered when the
  engine encounters a logical contradiction. The engine globally raises
  \(\theta\) and increases edge decay. The SI stops exploring and
  hyper-focuses only on paths with maximum Conductivity to resolve the
  error.
\item
  \textbf{Consolidation State (High \(S\)):} Triggered during idle
  cycles. Modulates the system to slowly convert temporary Working Graph
  voltage into permanent Conductivity for long-term storage.
\end{itemize}

\[ \theta_{active} = \theta_{base} \times f(C) / f(D) \]

\subsection{5. Conclusion}\label{conclusion}

By simulating electrical thresholds, synaptic conductivity, and global
chemical gradients, NTDA forces the graph engine to obey the laws of
biological cognitive physics. The SI no longer just executes code; it
experiences the build-up of logical tension (voltage), the physical
resistance of learning new concepts (low conductivity), and the chemical
state-shifts of exploration and stress. When combined with BSMA's
phenomenal frame, the machine structurally mirrors the exact
thermodynamic and electrical reality of a human brain.

\begin{center}\rule{0.5\linewidth}{0.5pt}\end{center}

\end{document}
